\section*{الفصل \ref{ch:b2:magma-emplacement} --- الصهارات وتموضعها}

\begin{solution}{exo:b2:magma-emplacement:1}
$T = \qty{1073}{K}$. من أجل $w = 1$، $\ln w = 0$: $\log_{10}\eta = -3.545 +
9601/(1073 - 195.7) = -3.545 + 10.94 = 7.40$. ومن أجل $w = 4$، $\ln 4 = 1.386$:
$A = -3.545 + 1.155 = -2.390$، $B = 9601 - 3283 = 6318$، $C = 195.7 + 44.7 =
240.4$، و$\log_{10}\eta = -2.390 + 6318/832.6 = -2.390 + 7.588 = 5.20$.
وتنخفض اللزوجة بعامل $10^{7.40 - 5.20} = 10^{2.20} \approx 160$.
\end{solution}

\begin{solution}{exo:b2:magma-emplacement:2}
الريوليت: $p_s = (6/0.44)^2 = \qty{186}{MPa}$، على عمق $\num{186e6}/(2600\times9.81)
= \qty{7.3}{km}$. البازلت: $p_s = (0.5/0.34)^2 = \qty{2.16}{MPa}$، على
$\num{2.16e6}/(2600\times9.81) = \qty{85}{m}$. فالريوليت يبدأ بتكوين الفقاعات
في عمق القشرة، والبازلت قرب السطح فقط (إن لم يحمل CO$_2$).
\end{solution}

\begin{solution}{exo:b2:magma-emplacement:3}
$(1150 + 50)/2 = \qty{600}{\celsius}$ و$(1150 + 400)/2 = \qty{775}{\celsius}$.
\omterm{def:b2:magma-emplacement:bodies}{والقاطع} الذي في الصخر البارد يُبرَّد بشدّة أكبر عند تماسه، فتكون له
\omterm{def:b2:magma-emplacement:aureole}{الحافة المبرّدة} الأعرض.
\end{solution}

\begin{solution}{exo:b2:magma-emplacement:4}
(أ) عتبة، موازية للتطبّق. (ب) قاطع، يقطع التطبّق.
(ج) باثوليث، مساحته أكثر من \qty{100}{km^2} ومتجمّع من اندساسات كثيرة.
(د) لاكوليث، قبّب سقفه. (ه) قاطع: ليس للغرانيت تطبّق،
فكل صفيحة فيه قاطع، أيًّا كان اتجاهها.
\end{solution}

\begin{solution}{exo:b2:magma-emplacement:5}
$(1 - 1.35\phi)^{-2.5}$: من أجل $\phi = 0.05$، $0.9325^{-2.5} = 1.19$، أي
\qty{119}{Pa.s}، مقابل $1 + 2.5\times0.05 = 1.125$، أي \qty{113}{Pa.s}، بصيغة
أينشتاين (وهما متقاربتان، كما ينبغي في حدّ التخفيف). ومن أجل $0.20$: $0.73^{-2.5} =
2.20$، أي \qty{220}{Pa.s}؛ ومن أجل $0.40$: $0.46^{-2.5} = 6.97$، أي \qty{697}{Pa.s}؛ ومن أجل $0.60$:
$0.19^{-2.5} = 63.6$، أي نحو \qty{6360}{Pa.s}. والانحشار عند $\phi = 1/1.35 = 0.74$.
\end{solution}

\begin{solution}{exo:b2:magma-emplacement:6}
$G = 300\times9.81 = \qty{2943}{Pa/m}$. (أ) $\bar u = 2943\times1^2/(12\times100) =
\qty{2.45}{m/s}$؛ $q = \qty{2.45}{m^2/s}$؛ $Re = 2700\times2.45\times1/100 = 66$، صفائحي؛
$t = 10000/2.45 = \qty{4080}{s}$، أي نحو \qty{1.1}{h}. (ب) $\bar u =
2943/(12\times\num{e6}) = \qty{2.45e-4}{m/s}$، $t = \qty{4.08e7}{s} = 472$ يومًا؛ ومع
$w = \qty{10}{m}$ تكون $\bar u$ أكبر بمئة مرة، \qty{0.0245}{m/s}، و$t =
\qty{4.08e5}{s} = 4.7$ أيام. (ج) يبرد \omterm{def:b2:magma-emplacement:bodies}{القاطع} الذي عرضه \qty{1}{m} في نحو
$0.5^2/\num{e-6} = \qty{2.5e5}{s}$، أي 2.9 يوم؛ \omterm{def:b2:magma-emplacement:bodies}{والقاطع} الذي عرضه \qty{10}{m} في $5^2/\num{e-6} =
\qty{2.5e7}{s}$، أي 289 يومًا. فالبازلت (\qty{1.1}{h}) والريوليت في \omterm{def:b2:magma-emplacement:bodies}{القاطع}
العريض (4.7 أيام) يصلان قبل أن يتجمّد قاطعاهما؛ أمّا الريوليت في
\omterm{def:b2:magma-emplacement:bodies}{القاطع} الذي عرضه \qty{1}{m} (472 يومًا) فلا يستطيع.
\end{solution}

\begin{solution}{exo:b2:magma-emplacement:7}
لكل كيلوغرام من \omterm{def:b2:magma-emplacement:magma}{الصهارة}، $n = 10/18.015 = \qty{0.555}{mol}$ من الماء و$nRT =
0.555\times8.314\times1173 = \qty{5414}{J}$؛ ويشغل المصهور الباقي $0.99/2400 =
\qty{4.13e-4}{m^3}$. عند \qty{100}{MPa}: $V = \qty{5.41e-5}{m^3}$، والكسر
$5.41/(5.41 + 41.3) = 0.12$. وعند \qty{10}{MPa}: $V = \qty{5.41e-4}{m^3}$، والكسر
$0.57$. وعند \qty{0.1}{MPa}: $V = \qty{0.0541}{m^3}$، والكسر $0.992$. والقيمة عند
\qty{100}{MPa} أقلها موثوقية: فبخار الماء عند مثل هذا الضغط كثيف
وبعيد عن المثالية (وكثير من الماء سيبقى في الواقع مذابًا).
\end{solution}

\begin{solution}{exo:b2:magma-emplacement:8}
(أ) $2300 + 60z = 2650$ تعطي $z = \qty{5.8}{km}$. (ب) كثافة المصهور
\qty{2580}{kg/m^3}: $z = 280/60 = \qty{4.7}{km}$، أي أعلى بنحو \qty{1.2}{km}. (ج)
عند $z = \qty{7.8}{km}$ تكون كثافة الصخر $2300 + 60\times7.8 = \qty{2770}{kg/m^3}$؛
والقوة $(2770 - 2650)\times9.81 = \qty{1.2}{kN}$ لكل متر مكعب، نحو الأعلى.
\end{solution}

\begin{solution}{exo:b2:magma-emplacement:9}
$v = 2\times300\times9.81\times(1000)^2/(9\times\num{e19}) = \qty{6.5e-11}{m/s}$، أي
$\num{6.5e-11}\times\num{3.156e7} = \qty{2.1e-3}{m}$ في السنة، \qty{2.1}{mm/yr}. والصعود
\qty{10}{km} يستغرق $\num{e4}/\num{2.1e-3} = \qty{4.8}{Myr}$. والسلّم الزمني للتبريد
هو $(1000)^2/\num{e-6} = \qty{e12}{s} = \qty{32}{kyr}$: فيتصلّب الجسم بعد
أن يصعد نحو $32000\times0.0021 \approx \qty{70}{m}$. لا يستطيع ديابير صهاري عبور
مثل هذه القشرة؛ إنه يحتاج إلى محيط أسخن وألين.
\end{solution}

\begin{solution}{exo:b2:magma-emplacement:10}
(أ) $x_{\mathrm{OH}} = \sqrt{K_2x_m} = \sqrt{K_1K_2f}$، ومنه $x = K_1f + \frac12\sqrt{K_1K_2f}$.
ونسبة الحدّ الأول إلى الثاني هي $2\sqrt{K_1f/K_2}$: تؤول إلى الصفر
حين $f \to 0$ (فيهيمن الجذر التربيعي) وتنمو بلا حدّ مع $f$ (فيهيمن
الحدّ الخطي). (ب) $K_1 f = \frac12\sqrt{K_1K_2 f}$ تعطي $K_1^2f^2 =
K_1K_2f/4$، ومنه $f = K_2/(4K_1)$. (ج) $f = 0.2/(\num{4e-4}) = \qty{500}{MPa}$. وعند
\qty{100}{MPa}: $x_m = 0.0100$، و$\frac12\sqrt{\num{e-4}\times0.2\times100} =
\frac12\times0.0447 = 0.0224$، ومنه $x = 0.0324$ وحصّة الهيدروكسيل
$0.0224/0.0324 = \qty{69}{\percent}$.
\end{solution}

\begin{solution}{exo:b2:magma-emplacement:11}
(أ) $1000/\num{4e6} = \qty{2.5e-4}{km^3}$ في السنة، أي $\qty{2.5e5}{m^3}$ في السنة،
أو $\num{2.5e5}/\num{3.156e7} = \qty{0.0079}{m^3/s}$، أي نحو ثمانية لترات في الثانية.
(ب) ينقل \omterm{def:b2:magma-emplacement:bodies}{القاطع} $1000\times2\times0.8 = \qty{1600}{m^3/s}$، ويوصل
$\qty{e12}{m^3}$ في \qty{6.25e8}{s}، أي نحو 20 سنة: $\num{5e-6}$ من
\qty{4}{Myr}. (ج) تحوي كل صفيحة $20\times\num{e4}\times\num{e4} = \qty{2e9}{m^3} =
\qty{2}{km^3}$، فيلزم 500 صفيحة، واحدة كل $\num{4e6}/500 = 8000$ سنة في المتوسط.
(د) $t = (10)^2/(4\times0.65^2\times\num{0.7e-6}) = \qty{8.5e7}{s}$، أي 2.7 سنة. فكل
صفيحة تتصلّب قبل وصول التالية بآلاف السنين: وما لم يأتِ الإمداد
في دفقات سريعة، فلن يكون \omterm{def:b2:magma-emplacement:bodies}{البلوتون} أبدًا جسمًا كبيرًا واحدًا من \omterm{def:b2:magma-emplacement:magma}{الصهارة}، كما
يوحي تأريخ البلوتونات الكبيرة.
\end{solution}

\begin{solution}{exo:b2:magma-emplacement:12}
(أ) عند $x$ ثابتة، $\ln(T - T_0) = \text{ثابت} - \frac12\ln t - x^2/(4\kappa t)$،
ومشتقته $-1/(2t) + x^2/(4\kappa t^2)$ تنعدم عند $t^* = x^2/(2\kappa)$. وعندئذ
$e^{-x^2/(4\kappa t^*)} = e^{-1/2}$ و$\sqrt{4\pi\kappa t^*} = x\sqrt{2\pi}$، ومنه $\Delta
T_{\max} = (T_m - T_0)\,b\,e^{-1/2}/(x\sqrt{2\pi}) = (T_m - T_0)\,b/(x\sqrt{2\pi e})$. (ب)
$\sqrt{2\pi e} = 4.13$: $x = 1100\times2/(100\times4.13) = \qty{5.3}{m}$، تُبلغ بعد
$t^* = 5.3^2/(2\times\num{0.7e-6}) = \qty{2.0e7}{s}$، أي نحو 230 يومًا (وعلى بعد 2.7
ضعف عرض \omterm{def:b2:magma-emplacement:bodies}{القاطع} فقط يكون تقريب المصدر المستوي خشنًا). (ج) يتناقص التسخين
بنسبة $b/x$: \omterm{def:b2:magma-emplacement:bodies}{فالقاطع} الرقيق يحمل حرارة قليلة، ولا يسخن الصخر سخونة ملموسة
إلا في حدود بضعة أضعاف عرض \omterm{def:b2:magma-emplacement:bodies}{القاطع}.
\end{solution}

\begin{solution}{pb:b2:magma-emplacement:1}
\textbf{1.} $\log_{10}\eta = -4.55 + 5510/(1423 - 611) = -4.55 + 6.786 = 2.24$؛
$\eta = \qty{172}{Pa.s}$.

\textbf{2.} $-4.55 + 5510/762 = -4.55 + 7.231 = 2.68$، $\eta = \qty{480}{Pa.s}$: أي
عامل $480/172 = 2.8$ على \qty{50}{K}.

\textbf{3.} $172\times(1 - 1.35\times0.20)^{-2.5} = 172\times0.73^{-2.5} = 172\times2.20 =
\qty{378}{Pa.s}$.

\textbf{4.} $\phi = 1/1.35 = 0.74$. وتضع الشواهد الحقلية نهاية الجريان الكلي
عند حصّة مصهور من \qtyrange{30}{50}{\percent}، أي
\qtyrange{50}{70}{\percent} من البلّورات: أبكر، لأن البلّورات الحقيقية ليست
كرات ذات حجم واحد.

\textbf{5.} القاعدة: $2800 - 2710 = \qty{90}{kg/m^3}$، \omterm{def:b2:magma-emplacement:magma}{فالصهارة} أخفّ.
الرسوبيات: $2450 - 2710 = \qty{-260}{kg/m^3}$، \omterm{def:b2:magma-emplacement:magma}{فالصهارة} أكثف.

\textbf{6.} $p = 2450\times9.81\times3000 = \qty{72.1}{MPa}$.

\textbf{7.} $90\times9.81 = \qty{883}{Pa/m}$.

\textbf{8.} $\bar u = 883\times2.0^2/(12\times378) = 3532/4536 = \qty{0.78}{m/s}$.

\textbf{9.} $q = 0.78\times2.0 = \qty{1.56}{m^2/s}$؛ وعلى \qty{1.0}{km}، نحو
\qty{1560}{m^3/s}.

\textbf{10.} $Re = 2710\times0.78\times2.0/378 = 11$: صفائحي.

\textbf{11.} $12000/0.78 = \qty{15400}{s}$، أي نحو \qty{4.3}{h}.

\textbf{12.} $(1.0)^2/\num{0.7e-6} = \qty{1.43e6}{s}$، أي 16.5 يومًا، أطول بكثير من
\qty{4.3}{h}: فيبقى \omterm{def:b2:magma-emplacement:bodies}{القاطع} مفتوحًا ما دامت \omterm{def:b2:magma-emplacement:magma}{الصهارة} تجري.

\textbf{13.} عند \qty{3.0}{km} تلاقي \omterm{def:b2:magma-emplacement:magma}{الصهارة} رسوبيات أخفّ منها: فقاعدة
الرسوبيات هي مستوى طفوها المتعادل. وفوقه تكون \omterm{def:b2:magma-emplacement:magma}{الصهارة}
أكثف مما حولها، فتنتشر جانبيًا على امتداد التماس على هيئة
عتبة.

\textbf{14.} $p_s = (0.40/0.34)^2 = \qty{1.38}{MPa}$، يُبلغ على عمق
$\num{1.38e6}/(2450\times9.81) = \qty{58}{m}$ تحت السطح.

\textbf{15.} $0.34\times\sqrt{72.1} = 0.34\times8.49 = \qty{2.9}{\percent}$، أكثر بكثير من
\qty{0.40}{\percent} الموجودة: فلا تفصل صهارة \omterm{def:b2:magma-emplacement:bodies}{العتبة} ماءً.

\textbf{16.} CO$_2$ أقل ذوبانًا من الماء بكثير: فهو يتشبّع ويكوّن فقاعات
في عمق مسار الصعود، ولذلك يكون أول غاز في البازلت أساسًا CO$_2$.

\textbf{17.} لكل كيلوغرام، $n = 4.0/18.015 = \qty{0.222}{mol}$، $V = 0.222\times8.314\times1423/
\num{e5} = \qty{0.0263}{m^3}$؛ ويشغل المصهور $0.996/2710 = \qty{3.68e-4}{m^3}$؛
والكسر $0.0263/(0.0263 + 0.000368) = 0.986$، أي \qty{98.6}{\percent}.

\textbf{18.} $(1150 + 100)/2 = \qty{625}{\celsius}$.

\textbf{19.} بعد 10 سنوات، $2\sqrt{\kappa t} = 2\sqrt{\num{0.7e-6}\times\num{3.156e8}} =
\qty{29.7}{m}$. والنقطة على $x = \qty{60}{m}$ من المنتصف: $T = 100 +
525\,[\operatorname{erf}(-10/29.7) + \operatorname{erf}(110/29.7)] = 100 + 525\,(-0.365 +
1.000) = \qty{433}{\celsius}$.

\textbf{20.} في المنتصف، $T = 100 + 1050\operatorname{erf}(50/2\sqrt{\kappa t})$: بعد
10 سنوات $1050\operatorname{erf}(1.68) = 1050\times0.983$، $T = \qty{1132}{\celsius}$؛ وبعد
50 سنة، $2\sqrt{\kappa t} = \qty{66.5}{m}$، $1050\operatorname{erf}(0.752) = 1050\times0.713$،
$T = \qty{848}{\celsius}$.

\textbf{21.} عند $x = \qty{130}{m}$: \qty{294}{\celsius} بعد 300 سنة،
و\qty{295}{\celsius} بعد 363 سنة (القيمة القصوى)، و\qty{295}{\celsius} بعد 420
سنة. فبإهمال الحرارة الكامنة لا يكاد الصخر الواقع على بعد \qty{80}{m} فوق العتبة
يبلغ \qty{300}{\celsius} أبدًا.

\textbf{22.} $\num{9e8}\times1.772/(\num{2.5e6}\times1050) = 0.608$.

\textbf{23.} $\lambda = 0.64$: $e^{-0.4096}/(0.64\times1.6346) = 0.6639/1.0461 = 0.635$؛
$0.65$: $0.6554/1.0673 = 0.614$؛ $0.66$: $0.6469/1.0886 = 0.594$. والقيمة $0.608$
تقع بين $0.65$ و$0.66$، قريبة من $0.65$: $\lambda = 0.65$.

\textbf{24.} $100 + 1050/(1 + 0.642) = 100 + 640 = \qty{740}{\celsius}$. فالحرارة
الكامنة المحرَّرة عند جبهة التصلّب تبقي جانب \omterm{def:b2:magma-emplacement:magma}{الصهارة} عند
\qty{1150}{\celsius} مدة أطول، فيتلقّى التماس حرارة أكثر.

\textbf{25.} $(50)^2/\num{0.7e-6} = \qty{3.57e9}{s}$، أي 113 سنة.

\textbf{26.} $t = (100)^2/(16\times0.65^2\times\num{0.7e-6}) = \num{e4}/\num{4.73e-6} =
\qty{2.11e9}{s}$، ومنه $\num{2.11e9}/\num{3.156e7}$: فتصير \omterm{def:b2:magma-emplacement:bodies}{العتبة} صلبة بعد
\textbf{\qty[text-series-to-math=true]{67}{} سنة}. وهو أقصر من
$(b/2)^2/\kappa$ بالعامل $1/(4\lambda^2) = 0.59$: فذلك السلّم الزمني لا يحدّد إلا
رتبة القيمة، وجدارا \omterm{def:b2:magma-emplacement:bodies}{العتبة} كلاهما يسحبان الحرارة إلى مضيف بارد في
آن واحد، بسرعة تكفي لإزالة الحرارة الكامنة أيضًا.
\end{solution}
